\documentclass[11pt]{article}
\usepackage[margin=1in]{geometry}
\usepackage{amsmath,amssymb}
\usepackage{booktabs}
\usepackage{longtable}
\usepackage{array}
\usepackage{enumitem}
\usepackage{hyperref}
\usepackage{xcolor}

\hypersetup{
    colorlinks=true,
    linkcolor=blue!60!black,
    urlcolor=blue!60!black
}

\title{\textbf{Phase 1 Report: Continuous Relaxation and Rounding for Graph Isomorphism}}
\author{}
\date{Codebase-grounded draft, 2026-05-01}

\begin{document}
\maketitle

\begin{abstract}
This report documents the current state of a graph isomorphism pipeline implemented in this repository.
The codebase, not earlier theory drafts, is treated as the source of truth.

The implemented system solves a continuous relaxation of the permutation search problem over the Birkhoff polytope,
uses a quadratic commutator penalty to reward structure-preserving correspondences,
adds a degree-profile compatibility term to sharpen the landscape,
and converts the relaxed solution into a permutation candidate through rounding.

The report has four goals.
First, formalize the implemented objective.
Second, justify why the relaxation is useful.
Third, compare adjacency-only matching against the current degree-augmented objective.
Fourth, position Hungarian rounding and hyperplane rounding within the present pipeline.

All numerical comparisons in this report were re-run from the current repository on 2026-05-01.
The main conclusion of Phase 1 is that the degree-profile term is already doing real work:
it dramatically improves separation on verified non-isomorphic random pairs while leaving solve times essentially unchanged.
The current active solver now uses hyperplane rounding rather than Hungarian rounding inside the main pipeline.
However, on the tested reruns, hyperplane rounding has not yet outperformed an offline Hungarian comparison baseline on the saved relaxed matrices.
\end{abstract}

\tableofcontents
\newpage

\section{Introduction}

Given two simple undirected graphs $G_1 = (V_1, E_1)$ and $G_2 = (V_2, E_2)$ with the same number of vertices,
the graph isomorphism problem asks whether there exists a relabeling of the vertices of $G_1$ that makes its edge set identical to that of $G_2$.

If $A$ and $B$ denote the adjacency matrices of the two graphs,
then an isomorphism can be represented by a permutation matrix $P$ satisfying
\begin{equation}
AP = PB
\end{equation}
or equivalently
\begin{equation}
P^\top A P = B.
\end{equation}

The discrete search space has size $n!$.
That is the central obstacle.

The approach implemented here replaces the discrete search over permutation matrices with a continuous optimization over doubly stochastic matrices.
The optimization returns a soft correspondence matrix $X^*$.
That matrix is then rounded into a permutation candidate and checked exactly in integer arithmetic.

This report focuses on the Phase 1 version of that idea.
It does not describe a hypothetical future solver.
It documents what is currently implemented, what the current comparisons show, and where the next additions naturally fit.

\section{Graph Isomorphism as a Permutation Problem}

Let $A, B \in \{0,1\}^{n \times n}$ be the adjacency matrices of two graphs on $n$ vertices.
For a permutation matrix $P \in \{0,1\}^{n \times n}$,
the identity $AP = PB$ means that column $j$ of $P$ selects the image of vertex $j$ under the candidate relabeling.

This formulation is attractive because it expresses isomorphism as a rigid algebraic condition.
If a correct $P$ exists, the graphs are isomorphic.
If no such $P$ exists, they are not.

The difficulty is combinatorial.
The feasible set of permutation matrices is finite but enormous.
Even moderate graph sizes rule out brute-force enumeration.

This motivates relaxation.
Instead of insisting that every row and every column of the decision matrix contain exactly one binary $1$,
we first allow the mass to spread continuously.
The hope is that the geometry of the relaxed problem still points toward useful permutations.

\section{Why Relaxation}

\subsection{Birkhoff Polytope}

A matrix $X \in \mathbb{R}^{n \times n}$ is doubly stochastic if
\begin{equation}
X_{ij} \ge 0, \qquad \sum_{j=1}^{n} X_{ij} = 1 \ \forall i, \qquad \sum_{i=1}^{n} X_{ij} = 1 \ \forall j.
\end{equation}

The set of all such matrices is the Birkhoff polytope $\mathcal{B}_n$.
Its vertices are exactly the permutation matrices.

That fact is the structural reason the relaxation is meaningful.
The continuous feasible region still contains the discrete solutions we ultimately care about.

\subsection{Interpretation}

The relaxed variable $X$ is best viewed as a soft correspondence matrix.
Row $i$ describes how strongly vertex $i$ in $G_1$ is matched to each vertex of $G_2$.
If $X$ lands on a permutation vertex, the matching is sharp.
If $X$ remains fractional, the solver has found ambiguity, symmetry, or a poor local explanation.

This interpretation is useful in two ways.
First, it turns the discrete search into a smooth numerical optimization problem.
Second, it gives an intermediate object that can be studied even when exact certification fails.

\section{Implemented Phase 1 Objective}

The current code in \texttt{isomorphism.py} implements two signals:
\begin{enumerate}[label=\arabic*.]
    \item a linear degree-profile compatibility term, and
    \item a quadratic adjacency commutator penalty.
\end{enumerate}

No higher-order walk term, spectral penalty, or entropy regularization is currently present in the active solver.
Those belong to future phases, not to the present source-of-truth formulation.

\subsection{Adjacency Commutator}

The structural core of the objective is
\begin{equation}
\mathcal{Q}(X) = \|AX - XB\|_F^2.
\end{equation}

This term is directly motivated by the exact permutation condition.
If $X = P$ is a correct permutation matrix, then $AP = PB$ and hence $\mathcal{Q}(P) = 0$.

That makes $\mathcal{Q}$ the natural baseline objective.
It is symmetric, concise, and tied directly to the algebra of isomorphism.

\subsection{Degree-Profile Compatibility}

The current solver also computes the degree vectors
\begin{equation}
d^A_i = \sum_k A_{ik}, \qquad d^B_j = \sum_k B_{jk},
\end{equation}
and the one-hop neighbor-degree summaries
\begin{equation}
s^A_i = \sum_k A_{ik} d^A_k, \qquad s^B_j = \sum_k B_{jk} d^B_k.
\end{equation}

These define the cost matrix
\begin{equation}
C_{ij} = \alpha |d^A_i - d^B_j| + \beta |s^A_i - s^B_j|,
\end{equation}
with current code-level weights
\begin{equation}
\alpha = 1.0, \qquad \beta = 0.35.
\end{equation}

The corresponding linear term is
\begin{equation}
\mathcal{L}(X) = \langle C, X \rangle_F = \sum_{i,j} C_{ij} X_{ij}.
\end{equation}

\subsection{Final Implemented Optimization Problem}

The current implemented Phase 1 objective is
\begin{equation}
\boxed{
\begin{aligned}
\min_{X \in \mathcal{B}_n} \quad & \mathcal{L}(X) + \mathcal{Q}(X) \\
= \quad & \sum_{i,j} C_{ij} X_{ij} + \|AX - XB\|_F^2.
\end{aligned}}
\end{equation}

When the degree weights are set to zero, the solver reduces to the adjacency-only baseline
\begin{equation}
\min_{X \in \mathcal{B}_n} \|AX - XB\|_F^2.
\end{equation}

That ablation is central to this report because it isolates the role of the extra cost term.

\section{Why the Objective Is Built This Way}

\subsection{Why Start from $\|AX - XB\|_F^2$}

The adjacency commutator is the cleanest continuous analog of the discrete isomorphism condition.
It asks the relaxed matrix $X$ to approximately intertwine the two adjacency operators.

That is exactly the right starting point.
If the report used any other baseline, the mathematical link to permutation-based isomorphism would be weaker.

\subsection{Why Add Degree Information}

The adjacency-only objective often leaves the relaxed landscape too flat.
Many soft assignments can produce small commutator residuals before the solver has committed to a sharp correspondence.

The degree-profile term addresses that weakness by injecting local compatibility before the quadratic term has fully resolved the global structure.
In practice, it behaves like a low-cost prior over plausible node matches.

The extra term is also cheap.
It changes neither the dimension of the feasible region nor the asymptotic structure of the QP.

That makes it an excellent Phase 1 addition:
it adds signal without forcing a redesign of the solver.

\section{From Relaxed Matrix to Isomorphism Index}

Once the solver finds an optimum $X^*$, it reports the objective value $Z^*$ and the repository converts it into a bounded score
\begin{equation}
I = \exp(-\lambda Z^*),
\end{equation}
with current scaling parameter
\begin{equation}
\lambda = 0.1.
\end{equation}

This score is not a proof of isomorphism.
It is a similarity indicator derived from the relaxed objective.

When $Z^*$ is tiny, $I$ is near $1$.
When $Z^*$ is large, $I$ drops away from $1$.

The point of the score is comparative interpretation.
It lets us see whether one objective spreads apart isomorphic and non-isomorphic cases more effectively than another.

\section{Rounding and Exact Verification}

\subsection{Current Active Path: Hyperplane Rounding}

The current active pipeline does not stop at $X^*$.
It rounds the doubly stochastic matrix to a permutation candidate using the hyperplane-rounding routine in \texttt{hyperplane\_rounding.py}.

The implemented procedure first computes an SVD-based embedding of the rows and columns of $X^*$.
It then samples random hyperplanes in the embedding space, converts the projections into sign signatures, and solves a linear assignment on the resulting Hamming-distance matrix.

The current \texttt{compute\_isomorphism\_index} path runs this routine with 200 trials and a fixed seed, then verifies the resulting permutation candidate exactly.

\subsection{Offline Comparison Baseline: Hungarian Rounding}

Although Hungarian rounding is no longer wired into the active solver path, it remains the natural comparison baseline.
Given $X^*$, Hungarian rounding computes the nearest permutation under linear assignment by minimizing the cost matrix $1 - X^*$.

That baseline is useful because it is deterministic, extremely fast, and easy to interpret.
It remains the cleanest reference point for evaluating whether hyperplane rounding is actually extracting extra geometric information from the relaxed matrix.

\subsection{Exact Integer Verification}

After rounding, the code verifies the candidate permutation exactly by checking whether
\begin{equation}
AP^* - P^*B = 0
\end{equation}
in integer arithmetic.

This matters.
The relaxed objective and the index $I$ are heuristic similarity measures.
The residual check is the exact certificate.

If the residual vanishes, the solver has produced a valid isomorphism witness.
If it does not, the result is inconclusive rather than a proof of non-isomorphism.

\subsection{Why Keep the Hungarian Comparison}

The switch to hyperplane rounding does not remove the need for comparison.
On the contrary, once hyperplane rounding becomes the active path, it becomes even more important to compare it against Hungarian rounding offline.

The evaluation question is simple:
does the geometry-aware rounding path certify more true isomorphisms, or recover better permutations on hard fractional cases, than the nearest-permutation baseline?

Phase 1 can now answer that question empirically on the current benchmark family.

\section{Experimental Setup for This Report}

All results below were generated from the current repository on 2026-05-01.
The runs used the code in \texttt{generators.py}, \texttt{isomorphism.py}, and \texttt{hyperplane\_rounding.py}.

Two objective configurations were compared:
\begin{enumerate}[label=\arabic*.]
    \item \textbf{Adjacency only}: $\alpha = 0$, $\beta = 0$.
    \item \textbf{Degree plus adjacency}: $\alpha = 1.0$, $\beta = 0.35$.
\end{enumerate}

For each $n \in \{5,6,7,8,9,10\}$,
I generated six isomorphic pairs and six nominally non-isomorphic pairs.
Every generated pair was then checked with NetworkX isomorphism testing so that the report could distinguish truly non-isomorphic samples from accidental collisions.

This exact-label check is important.
The current \texttt{make\_non\_isomorphic\_pair} function samples two dense random graphs independently.
That procedure makes non-isomorphism overwhelmingly likely, but not guaranteed.

\section{Ablation I: Adjacency-Only vs Degree-Augmented Objective}

\subsection{Overall Aggregate Results}

Table~\ref{tab:overall-objective} compares the two objectives over the rerun dataset.
The most important line is the verified non-isomorphic subset:
that excludes accidental isomorphic collisions from the nominal non-isomorphic bucket.

\begin{table}[h]
\centering
\caption{Overall objective comparison on fresh reruns under the active hyperplane-rounding pipeline. Times are mean end-to-end solve times.}
\label{tab:overall-objective}
\begin{tabular}{>{\raggedright\arraybackslash}p{4.0cm} >{\raggedright\arraybackslash}p{3.0cm} r r r r}
\toprule
Objective & Dataset & Mean time (s) & Mean $Z^*$ & Mean $I$ & Active cert rate \\
\midrule
Adjacency only & Isomorphic & 0.00973 & $6.38 \times 10^{-9}$ & 0.99999999936 & 77.8\% \\
Adjacency only & Verified non-isomorphic & 0.01197 & 0.51625 & 0.94989 & 0.0\% \\
\addlinespace
Degree + adjacency & Isomorphic & 0.00896 & $3.75 \times 10^{-9}$ & 0.99999999963 & 77.8\% \\
Degree + adjacency & Verified non-isomorphic & 0.01109 & 11.05194 & 0.37353 & 0.0\% \\
\bottomrule
\end{tabular}
\end{table}

The contrast is sharp.
On isomorphic pairs, both objectives behave almost identically:
solve times are essentially the same,
the objective values collapse toward zero,
and the index remains effectively $1$.

On verified non-isomorphic pairs, the difference is dramatic.
Adjacency-only matching leaves the mean index at approximately $0.95$,
which is far too close to the isomorphic regime to be a useful similarity score.

Once the degree-profile term is restored, the mean index on verified non-isomorphic pairs drops to approximately $0.37$.
That is the strongest empirical justification for the extra cost term in the current code.

\subsection{Per-$n$ Separation Trend}

Table~\ref{tab:per-n-noniso} shows the same pattern broken down by graph size for verified non-isomorphic pairs.
The degree-augmented objective produces much stronger separation at every tested $n$.

\begin{table}[h]
\centering
\caption{Verified non-isomorphic pairs only: per-$n$ mean scores and solve times.}
\label{tab:per-n-noniso}
\begin{tabular}{r r r r r r}
\toprule
$n$ & \multicolumn{2}{c}{Mean $I$} & \multicolumn{2}{c}{Mean $Z^*$} & Mean time (s) \\
\cmidrule(lr){2-3}\cmidrule(lr){4-5}\cmidrule(lr){6-6}
 & Adj. only & Degree + adj. & Adj. only & Degree + adj. & Both objectives \\
\midrule
5  & 0.9642 & 0.6392 & 0.3646 & 4.4750 & $\approx 0.0043$ \\
6  & 0.9704 & 0.5257 & 0.3004 & 6.5277 & $\approx 0.0056$ \\
7  & 0.9641 & 0.4100 & 0.3655 & 9.0144 & $\approx 0.0078$ \\
8  & 0.9446 & 0.3135 & 0.5706 & 12.1260 & $\approx 0.0112$ \\
9  & 0.9145 & 0.1888 & 0.8973 & 17.1273 & $\approx 0.0159$ \\
10 & 0.9570 & 0.3140 & 0.4409 & 11.7582 & $\approx 0.0237$ \\
\bottomrule
\end{tabular}
\end{table}

This table answers the central Phase 1 question.
Why is the objective not just $\|AX - XB\|_F^2$?

Because that term alone does not spread apart the random verified non-isomorphic pairs enough.
The degree-profile term supplies low-order structural bias that the commutator alone is not extracting reliably in the relaxed space.

\subsection{Interpretation}

These results do not mean that degree information solves graph isomorphism.
They mean something narrower and more important for the current phase:
the extra linear term improves the geometry of the practical optimization landscape without increasing solve time in any meaningful way.

That is exactly the kind of modification Phase 1 should prioritize.

\section{Ablation II: Why Rounding Matters at All}

The aggregate tables above already show that even for exact isomorphic pairs,
the certification rate is only 77.8\% on the tested reruns.
That is not because the graphs are failing to be isomorphic.
It is because the relaxed optimum $X^*$ is not always close enough to a certifying permutation for the current rounding path to recover it.

This is easiest to see on symmetric graph families.

\subsection{Cycle Graph Case Study}

I tested an 8-cycle against an isomorphic relabeling.
For both objective variants, the solver returned a completely uniform matrix:
\begin{equation}
X^*_{ij} = \frac{1}{8} \qquad \text{for all } i,j.
\end{equation}

The first row of $X^*$ was
\begin{equation}
[0.125, 0.125, 0.125, 0.125, 0.125, 0.125, 0.125, 0.125].
\end{equation}

The measured objective value was numerically zero:
\begin{equation}
Z^* = 1.79 \times 10^{-30},
\end{equation}
so the index was exactly
\begin{equation}
I = 1.0.
\end{equation}

Yet neither rounding procedure produced a certificate.
The active hyperplane path failed, with best residual 32 after 200 trials,
and the offline Hungarian comparison also returned \texttt{false}.

This example is useful because it isolates the issue cleanly.
The optimization objective is satisfied,
but the solution lies at a highly symmetric interior point of the Birkhoff polytope rather than at a permutation vertex.

That is why rounding is not a cosmetic post-processing step.
It is a central part of the overall method.

\section{Rounding Comparison: Active Hyperplane vs Offline Hungarian}

\subsection{Aggregate Comparison}

Table~\ref{tab:rounding} compares the two rounding schemes on the rerun matrices.

\begin{table}[h]
\centering
\caption{Rounding comparison on the rerun dataset. Hyperplane is the active in-solver path; Hungarian is computed offline on the saved $X^*$ matrices for comparison.}
\label{tab:rounding}
\begin{tabular}{>{\raggedright\arraybackslash}p{4.0cm} r r r r r}
\toprule
Objective source & Cases & Hyperplane certs & Hungarian certs & Hyperplane-only wins & Hungarian round time \\
\midrule
Adjacency-only matrices & 72 & 30 & 30 & 0 & $3.14 \times 10^{-6}$ s \\
\addlinespace
Degree + adjacency matrices & 72 & 30 & 30 & 0 & $2.78 \times 10^{-6}$ s \\
\bottomrule
\end{tabular}
\end{table}

On the tested dataset, hyperplane rounding matched Hungarian exactly in certification count.
It produced no hyperplane-only wins and no Hungarian-only wins.

The timing contrast is still clear.
The offline Hungarian pass takes only a few microseconds per matrix at these sizes.
The active hyperplane cost is already included inside the end-to-end solve times reported earlier, which now average about 0.0100--0.0108 seconds per case across the two objective variants on this rerun.

\subsection{Interpretation}

This does \emph{not} imply that hyperplane rounding is useless.
It implies something more specific:
the current implementation and the currently tested matrices have not yet demonstrated a practical advantage over the existing baseline.

That is still useful progress.
It means the report can describe the new active hyperplane path honestly, together with clear evaluation criteria:
\begin{enumerate}[label=\arabic*.]
    \item more certification wins than Hungarian on hard fractional cases,
    \item meaningful advantage on structured benchmark families,
    \item acceptable runtime overhead relative to the gain.
\end{enumerate}

At the moment, those wins have not yet appeared in the tested runs.

\section{Threats to Validity}

\subsection{Random ``Non-Isomorphic'' Pairs Are Only Probabilistic}

The generator used for non-isomorphic pairs samples two dense random graphs independently.
That procedure does not enforce non-isomorphism.

In a fresh six-pair rerun at $n=5$,
two of the six nominally non-isomorphic samples were in fact isomorphic under exact NetworkX checking.
For $n=6$ through $n=10$, the rerun happened to contain no such collisions.

This is not a cosmetic detail.
If accidental collisions are not filtered out, they distort the very score separation that the report is trying to measure.

That is why the verified non-isomorphic subset is the right comparison set in this draft.

\subsection{Current Benchmark Family Is Narrow}

The current code mainly exercises dense Erd\H{o}s--R\'enyi style random graphs.
That is convenient for early benchmarking but not sufficient as a stress suite.

Highly symmetric families, regular families, and specially constructed hard pairs can behave very differently from random dense graphs.
The cycle-graph example above already shows that the relaxed solver can settle on a perfectly diffuse matrix in the presence of symmetry.

\subsection{The Report Must Follow the Code, Not Earlier Notes}

Earlier repository documents discuss additional terms such as higher-order walk penalties or spectral alignment.
Those are not part of the current active solver.

This report therefore treats the code as authoritative.
That choice is deliberate.
A progress report is strongest when every mathematical claim can be traced to running code.

\section{What Phase 1 Establishes}

Phase 1 establishes five concrete points.

\begin{enumerate}[label=\arabic*.]
    \item The graph isomorphism problem can be attacked as a continuous optimization over doubly stochastic matrices.
    \item The adjacency commutator is the correct structural baseline because it is the direct relaxation of $AP = PB$.
    \item In practice, that baseline alone is too weak as a similarity score on verified non-isomorphic random pairs.
    \item The implemented degree-profile term materially improves separation while leaving end-to-end solve times almost unchanged.
    \item Rounding remains the main bottleneck between a good soft matrix and a certifying permutation, and the current active hyperplane path has not yet beaten the offline Hungarian baseline on the tested runs.
\end{enumerate}

\section{Clean Insertion Points for the Next Additions}

This report is intentionally structured so that future phases can be inserted without rewriting the core story.

The natural extensions are:
\begin{enumerate}[label=\arabic*.]
    \item stronger objective terms, added as new ablation chapters after Section 8,
    \item better benchmark families, inserted into the experimental setup and validity sections,
    \item improved rounding methods, added to the rounding comparison chapter,
    \item exact or semi-exact constraints that sharpen the relaxed solution before rounding.
\end{enumerate}

That modularity matters.
The current report should remain valid even after the next terms are added.
Only the comparison chapters need to expand.

\section{Conclusion}

The present codebase already supports a coherent Phase 1 report.
The method is not merely a relaxation plus a vague score.
It is a complete pipeline:
formulate the correspondence problem,
relax it to the Birkhoff polytope,
solve a structured objective,
round the soft solution,
and verify the candidate permutation exactly.

The key empirical result of this phase is simple and defensible.
Adding the degree-profile cost term sharply improves the separation between isomorphic and verified non-isomorphic random graph pairs while keeping runtime effectively unchanged.

The key open issue is equally clear.
A good relaxed matrix is not automatically a good permutation witness.
That is why rounding remains a first-class research question rather than a final afterthought.

In that sense, switching the main pipeline to hyperplane rounding is a meaningful Phase 1 update.
It promotes rounding from a side experiment into part of the source-of-truth method.
The current evidence simply says that this promotion has not yet changed certification outcomes relative to the offline Hungarian baseline on the tested data.

That is an honest and useful Phase 1 result.
\end{document}
